% ----- CHAUPAI 28 -----

\newpage

\subsection*{\deva{अष्टाविंशतितम चौपाई} (Twenty-Eighth Chaupai)}
\addcontentsline{toc}{subsection}{\deva{अष्टाविंशतितम चौपाई} (Twenty-Eighth Chaupai) — \deva{और मनोरथ...}}

\begin{minipage}[t]{0.55\textwidth}
\vspace{0pt}

{\large \linespread{1.5}\selectfont
\deva{और मनोरथ जो कोइ लावै।}\\
\deva{सोइ अमित जीवन फल पावै॥}
\par}

\vspace{0.5em}

{\large \linespread{1.5}\selectfont
\telu{ఔర మనోరథ జో కోఇ లావై।}\\
\telu{సోఇ అమిత జీవన ఫల పావై॥}
\par}

\vspace{0.5em}

{\large \linespread{1.3}\selectfont
\textit{aura manoratha jo koi lāvai |}\\
\textit{soi amita jīvana phala pāvai ||}
\par}
\end{minipage}%
\hfill
\begin{minipage}[t]{0.42\textwidth}
\vspace{0pt}
\begin{tcolorbox}[anchorbox]
\textbf{Meaningful Ambitions:} True success isn't just about wishing for petty material things; it is about developing the clarity to seek the \textit{Amit Jeevan Phal}—the limitless fruit of a deeply purposeful, well-lived, and meaningful life.
\vspace{0.5em}\newline He accepted the ultimate reward only because his truest ambition was meaningful dedication, not material wealth.\newline\null\hfill\href{https://www.valmiki.iitk.ac.in/sloka?field_kanda_tid=6&language=dv&field_sarga_value=128}{\textit{Yuddha Kanda, 128:82-83}}
\end{tcolorbox}
\end{minipage}

\vspace{0.5em}
\subsubsection*{\textbf{Visual Rhythm Map:}}
\begin{table}[h!]
\centering
\renewcommand{\arraystretch}{1.2}
\noindent\makebox[\textwidth][l]{%
  {%
  \setlength{\tabcolsep}{0pt}%
  \begin{tabularx}{1.125\textwidth}{|*{18}{>{\centering\arraybackslash\hsize=1.0\hsize}X|}}
  \hline
  \multicolumn{3}{|>{\centering\arraybackslash\hsize=3.0\hsize}X|}{\textbf{\guru{औ}\laghu{र}} \par (3)} &
  \multicolumn{5}{>{\centering\arraybackslash\hsize=5.0\hsize}X|}{\textbf{\laghu{म}\guru{नो}\laghu{र}\laghu{थ}} \par (5)} &
  \multicolumn{2}{>{\centering\arraybackslash\hsize=2.0\hsize}X|}{\textbf{\guru{जो}} \par (2)} &
  \multicolumn{3}{>{\centering\arraybackslash\hsize=3.0\hsize}X|}{\textbf{\guru{को}\laghu{इ}} \par (3)} &
  \multicolumn{4}{>{\centering\arraybackslash\hsize=4.0\hsize}X|}{\textbf{\guru{ला}\guru{वै}} \par (4)} \\
  \hline
  \end{tabularx}%
  }%
  \hspace{0.01\textwidth}%
  \begin{minipage}{0.1\textwidth}
  \vspace{-0.5em}
  \centering
  {\Huge \textbf{17}}
  \end{minipage}%
}
\vspace{-1.5em}
\end{table}

\begin{table}[h!]
\centering
\renewcommand{\arraystretch}{1.2}
\setlength{\tabcolsep}{0pt}
\begin{tabularx}{\textwidth}{|*{16}{>{\centering\arraybackslash\hsize=1.0\hsize}X|}}
\hline
\multicolumn{3}{|>{\centering\arraybackslash\hsize=3.0\hsize}X|}{\textbf{\guru{सो}\laghu{इ}} \par (3)} &
\multicolumn{3}{>{\centering\arraybackslash\hsize=3.0\hsize}X|}{\textbf{\laghu{अ}\laghu{मि}\laghu{त}} \par (3)} &
\multicolumn{4}{>{\centering\arraybackslash\hsize=4.0\hsize}X|}{\textbf{\guru{जी}\laghu{व}\laghu{न}} \par (4)} &
\multicolumn{2}{>{\centering\arraybackslash\hsize=2.0\hsize}X|}{\textbf{\laghu{फ}\laghu{ल}} \par (2)} &
\multicolumn{4}{>{\centering\arraybackslash\hsize=4.0\hsize}X|}{\textbf{\guru{पा}\guru{वै}} \par (4)} \\
\hline
\end{tabularx}
\vspace{-1.5em}
\end{table}

\vspace{1em}
\textbf{ \deva{सरल-संस्कृतम्}:}\\
 \deva{यः कश्चित् अन्यम् अपि मनोरथम् (इच्छां) आनयति}\\
 \deva{सः जीवने अमितम् (असीमं) फलं प्राप्नोति॥}\\


Whoever brings any other desire or wish to you\\
that person obtains limitless fruits in their life.


\vspace{1em}
\newpage
\textbf{\deva{प्रतिपदार्थः} (Meaning of each word):}
\begin{table}[h!]
\centering
\renewcommand{\arraystretch}{1.2}
\begin{tabularx}{\textwidth}{@{} l l X X @{}}
\toprule
\textbf{Awadhi} & \textbf{Sanskrit} & \textbf{English} & \textbf{Grammar \& Notes} \\
\midrule
\deva{और} / \telu{ఔర}  &  \deva{अन्य / अपरः}  &  And / Other  &   \\
\deva{मनोरथ} / \telu{మనోరథ}  &  \deva{मनोरथम्}  &  Desire / Wish  &   \\
\deva{जो कोइ} / \telu{జో కోఇ}  &  \deva{यः कश्चित्}  &  Whoever  &   \\
\deva{लावै} / \telu{లావై}  &  \deva{आनयति}  &  Brings  &   \\
\deva{सोइ} / \telu{సోఇ}  &  \deva{सः एव}  &  That very person  & Awadhi pronoun \\
\deva{अमित} / \telu{అమిత}  &  \deva{अमितम् / असीमम्}  &  Limitless / Infinite  &   \\
\deva{जीवन फल} / \telu{జీవన ఫల}  &  \deva{जीवन-फलम्}  &  Fruit of life  &   \\
\deva{पावै} / \telu{పావై}  &  \deva{प्राप्नोति}  &  Obtains  &   \\
\bottomrule
\end{tabularx}
\end{table}

\begin{tcolorbox}[poetrybox]
\textbf{The 18-Beat Flow \& Bhava-Pradhan:}\\
Line 1 mathematically adds up to 17 beats. However, keeping in mind that this is a \textit{Bhava Pradhan} poem, the 16-matra rule is a convenience, not a strict law like in Sanskrit prosody. 
Singers joyfully and instinctively compress \deva{कोई} into a quick \textit{koy} without interrupting the singing flow, 
making this inconsistency a beautiful mnemonic quirk of the oral tradition!
\end{tcolorbox}

\newpage
